%% kapa: exam=2 round=25 year=2014 period=3 subject=compensation_law subject_ko=감정평가 및 보상법규 seq=2 number=2 points=30
%% source: Q-Net 감정평가사 기출문제 (cst00309) artlSeq=5109209 file=extracted/1008/5109209/2014년 제25회 감정평가사 2차시험 기출문제/3교시(감정평가 및 보상법규).hwp
%% conversion: kordoc 4.24.51 -> tools/gen2.py (auto); review=proofread-sonnet; flags=none
\begin{longquestion}{【문제 2】}
甲은 A시의 시외로 나가는 일반도로에 접한 자신 소유의 X토지에 교통로를 개설하고 대형음식점을 운영하고 있다. A시에서는 X토지와 이에 접하여 연결된 YㆍW토지의 소유권을 취득하여 혼잡한 교통량을 분산할 목적으로「국토의 계획 및 이용에 관한 법률」에 의거하여 우회도로를 설치한다는 방침을 결정하고, A시의 시장은 XㆍYㆍW토지의 개별공시지가 및 이 개별공시지가 산정의 기초가 된 P토지의 표준지공시지가와 도매물가상승률 등을 반영하여 산정한 보상기준가격을 내부적으로 결정하고 예산확보를 위해 중앙부처와 협의 중이다. 다음 물음에 답하시오. (30점) \par
1) 甲은 보상이 있을 것을 예상하여 더 많은 보상금을 받기 위해「부동산가격공시 및 감정평가에 관한 법률」에 의거하여 감정평가사를 통해 산정된 P토지의 표준지공시지가에 불복하여 취소소송을 제기하려고 한다. 이 경우 甲에게 법률상 이익이 있는지 여부를 검토하시오. (15점) \par
2) 위 취소소송에 P토지의 소유자인 丙이 소송에 참가할 수 있는지 여부와 甲이 확정 인용판결을 받았다면 이 판결의 효력은 YㆍW토지의 소유자인 乙에게도 미치는지에 대하여 설명하시오. (15점) \par
\end{longquestion}
